\documentclass[12pt]{extarticle}

% ===== Encodage & langue =====
\usepackage{fontspec}
\usepackage{polyglossia}
\setdefaultlanguage{french}

% ===== Police Verdana (compiler avec XeLaTeX ou LuaLaTeX) =====
\setsansfont{Verdana}
\renewcommand{\familydefault}{\sfdefault}

% ===== Interligne =====
\usepackage{setspace}
\setstretch{1.1}

% ===== Marges réduites =====
\usepackage{geometry}
\geometry{a4paper, top=1.8cm, bottom=1.8cm, left=2cm, right=2cm}

% ===== Mathématiques =====
\usepackage{amsmath}

% ===== Espacement réduit autour des équations =====
\setlength{\abovedisplayskip}{4pt}
\setlength{\belowdisplayskip}{4pt}
\setlength{\abovedisplayshortskip}{2pt}
\setlength{\belowdisplayshortskip}{2pt}

% ===== Espacement réduit des sections =====
\usepackage{titlesec}
\titlespacing*{\section}{0pt}{1.2ex}{0.6ex}

% ===== Titre du document =====
\title{\textbf{Correction du contrôle séquence 1 version 2}}
\author{}
\date{}

\begin{document}
\maketitle

% --------------------------------------------------
\section*{Exercice 1}
\begin{align*}
17 - 4 - 5 + 23 - 2 &= 13 - 5 + 23 - 2 \\
                     &= 8 + 23 - 2 \\
                     &= 31 - 2 \\
                     &= 29
\end{align*}

% --------------------------------------------------
\section*{Exercice 2}

\begin{align*}
\textbf{a)}
8 + 4 \times 3 &= 8 + 12 \\
               &= 20
\end{align*}

\textbf{b)}
\begin{align*}
41 - 7 \times 3 &= 41 - 21 \\
                &= 20
\end{align*}

\textbf{c)}
\begin{align*}
37 + 15 \div 5 &= 37 + 3 \\
               &= 40
\end{align*}

\textbf{d)}
\begin{align*}
5 + 8 \times 4 + 3 - 6 \div 3 &= 5 + 32 + 3 - 2 \\
                              &= 38
\end{align*}

% --------------------------------------------------
\section*{Exercice 3}
\begin{align*}
4 \times 6 + 2 \times 9 &= 24 + 18 \\
                        &= 42
\end{align*}

La dépense est de 42~€.

% --------------------------------------------------
\section*{Exercice 4}

\textbf{a)} $32 - (5 + 9) \div 2$
\begin{align*}
&= 32 - 14 \div 2 \\
&= 32 - 7 \\
&= 25
\end{align*}

\textbf{b)} $(4 + 7) \times 4 - 9 + 3$
\begin{align*}
&= 11 \times 4 - 9 + 3 \\
&= 44 - 9 + 3 \\
&= 38
\end{align*}

\textbf{c)} $(21 - (8 + 5)) \times 4$
\begin{align*}
&= (21 - 13) \times 4 \\
&= 8 \times 4 \\
&= 32
\end{align*}

\textbf{d)} $5 \times (25 - (9 + 11)) - 7$
\begin{align*}
&= 5 \times (25 - 20) - 7 \\
&= 5 \times 5 - 7 \\
&= 25 - 7 \\
&= 18
\end{align*}

% --------------------------------------------------
\section*{Exercice 5}
\begin{align*}
14 - 5 \times 2 &= 14 - 10 \\
                &= 4
\end{align*}

Une gomme coûte 4~€.

% --------------------------------------------------
\section*{Exercice bonus 1}

\textbf{a)} $8 \times 4 + 11$

\textbf{b)} $7 \times (9 + 2)$

% --------------------------------------------------
\section*{Exercice bonus 2}

\textbf{a)} La somme du produit de 9 par 4 et de 6.

\textbf{b)} Le produit de la somme de 7 et de 4 par 5.

\end{document}